%% notes-tex/019-simb-multimodal/sections/8-directions.tex
%% SOURCE: hand-authored. This section proposes work; it reports no measurement that is not
%% cited to an earlier section.

\section{What to run next\secstatus{todo}}
\label{sec:directions}

Ordered by cost against what each would settle. The first four are cheap enough that not
running them is the expensive choice.

\subsection{Tier 0, cheap and decisive\secstatus{todo}}

\begin{enumerate}
\item \textbf{The graph-prior probe} (\cref{sec:common}). Minutes of CPU, no GPU, no
  training. It asks whether deleting gene $X$ actually perturbs $X$'s neighbors on each of
  the nine graphs, which is the assumption the entire graph channel rests on and which has
  never been checked. Nine attention heads are currently spent enforcing it. If three
  graphs carry signal, regularize those three and free the other six.

\item \textbf{Morphology at full scale.} Drop \file{require_modalities} for the
  morphology-only arm and train on all 4,718 Ohya deletions instead of the 1,440 that also
  carry expression. This is a config change, and it is the only way the $0.0824$ becomes a
  number about morphology rather than about a quarter of morphology. Run it long enough to
  see a peak, which \cref{sec:common} says cannot be assumed from the expression budget.

\item \textbf{The scalar morphology warm-up}, on \file{A113} (ceiling $0.873$, robust CV
  $2.37$) or \file{C124_C} (Medium bud ratio, $0.698$ and $0.886$).
  Cheap, and it calibrates the epoch budget for the 278-feature target. Not colony size:
  colony size is fitness, which is already the project's strongest strand, and the CalMorph
  size features are the ones deletions barely move.

\item \textbf{Apply the post-hoc rescale to the existing best checkpoint.} Multiplying
  predictions by $r/s$ changes no correlation and moves \file{nmse} from $1.010$ to
  $0.944$, taking the model from worse-than-the-mean to better-than-the-mean in squared
  error (\cref{sec:objective}). An evening's work, no training, and it removes an objection
  the figure would otherwise have to answer.

\item \textbf{Bracket the expression peak, by resuming rather than restarting.} No run has
  observed a maximum at any budget up to 10,000 epochs, and a fresh 10,000-epoch run costs
  3.8 days. Resuming the existing best checkpoint is the cheap form of the same experiment,
  and the epoch at which the curve turns is what sets every arm budget afterward.
  \cref{sec:campaign} proposes $E = 4{,}000$ in the meantime, justified by an arm that has
  been observed to peak at 3,921.
\end{enumerate}

\subsection{Tier 1, the two questions worth GPU\secstatus{todo}}

Both are laid out with their epoch budgets, seed counts and slot costs in
\cref{sec:campaign}, which is the operational version of this subsection. In summary:

\begin{itemize}
\item \textbf{Settle the objective first} (\cref{tab:phase-a}). The quantile loss reaches
  its minimum at epoch 463 and rises for the next 9,500 while Pearson climbs, and the model
  never gets \file{nmse} below 1. A mechanism arm evaluated under an objective that is
  optimizing something else cannot separate ``the mechanism does not help'' from ``the
  objective does not reward it''.
\item \textbf{Then the pair term} (\cref{tab:phase-b}). The reference operator provably
  cannot express a dependence on the pair (deleted gene, reporter gene), and masked
  unmasking does not supply one at the point the score is taken. A rank-64 pair term
  already reaches $0.2274$ in 2.0 days against the masked objective's $0.2362$ in 3.8.
\item \textbf{Pair morphology with fitness for the first full-scale joint test}, since
  fitness shares 4,220 strains with morphology against expression's 1,440. Then run the
  expression pairing on the 1,440 as the mechanism test it was designed to be.
\end{itemize}

Substituting tyrosine for the 19 amino acids, as \file{igb_bx_aa.yaml} proposes, is now
predicted to be the \emph{null} arm rather than the live one, since tyrosine alone predicts
betaxanthin at $0.064$ against the profile's $0.298$. Keep it as a negative control, not as
the headline contrast.

\subsection{Tier 2, the structural options\secstatus{todo}}

\notesec{Constraint-based coupling is the mechanism Figure 6 needs} The 023 result is that
the coupling between amino-acid pools and betaxanthin is present in the data at $r = 0.298$
and that a plain auxiliary head does not use it, costing $-0.0265 \pm 0.0159$
(\cref{tab:bx-aa-paired}). That is a structural gap, and an unstructured second head is the
wrong instrument to close it: it asks the encoder to discover a stoichiometric relationship
from 19 noisy concentrations and one pigment readout, with nothing in the architecture
saying the two are connected through a pathway.

A constraint-based layer says it. Penalizing predictions that violate a stoichiometric or
flux balance over yeast-GEM makes the tyrosine-to-betalamic-acid route a property of the
model rather than something it must infer, and it turns the metabolome head from an
auxiliary target into a set of intermediates the constraint actually references.
\textbf{This is the Figure 6 mechanism, and it is what makes the amino-acid panel a step in
an argument instead of a negative result.}

Two facts size it rather than block it. yeast-GEM covers $19.4\,\%$ of the deletions in
the screen and $27\,\%$ of its high producers (\cref{tab:gem-coverage}), so the constraint
is silent on most of the genome and must be a term that switches off outside the model
rather than a hard requirement. And betaxanthin is already the strand closest to its
ceiling, so the headroom the constraint is competing for is the smallest of any strand:
$0.430$ against $0.914$. Expect it to earn its place by making the mechanism legible and by
extending to targets with no screen at all, rather than by moving betaxanthin's number a
lot.

\notesec{Not an expression rescue} The same layer is attractive for expression on the
argument that metabolism constrains transcription. \cref{tab:label-accounting} is the
reason to be cautious: expression's binding constraint is 1,484 perturbations each seen
once, and a metabolic prior does not add perturbations. Gate it on the same evidence
standard the amino-acid coupling met, a measurable relationship in the data first.

\notesec{Reifying the gene node into mRNA and protein} Not yet. The measured case for it is
thin: per-gene mRNA-to-protein correlation across strains has a median of $0.08$, a ridge
map between the two recovers $2$ to $3.5\,\%$ of held-out variance, and the two datasets
were measured in different media, so the decoupling that would justify separate nodes is
partly a condition artifact. Splitting the node multiplies parameters on an axis the data
does not yet resolve. The case changes when a same-media proteome and expression pair
exists, or when perturbations become fine-grained enough that the central-dogma collapse is
what breaks.

\notesec{Masked-label imputation as its own deliverable} The masked-conditioning oracle
reaches cross-study $0.4838$ at $m = 1000$ against the model's $0.238$, and $97.5$ to
$100.6\,\%$ of that gain survives removing a genotype predictor. It is a real capability
and a defensible product surface, and it is not a genotype-to-expression result. Deciding
whether to ship it as a separate thing is a decision, not an experiment.

\subsection{What to do about Figure 3 and Figure 6\secstatus{todo}}

\notesec{Figure 3} The manuscript currently claims $r = 0.543$ for expression and $0.619$
for morphology inside a paragraph marked \texttt{[FILLER -- R3.]}. The measured values are
$0.238$ and $0.082$. Neither placeholder should survive into a draft anyone reads. The
honest version of the panel is the ceiling-relative one: expression at $31\,\%$ of a
measured $0.775$, morphology at $13\,\%$ of a measured $0.611$ \emph{on a quarter of its
data}, with the fitness panel carrying the strong result. \textbf{A $0.24$ per-gene Pearson
is not a Figure 3 headline.}

Two of the three cheapest routes to a better one are not modeling work at all. Morphology
has never seen more than 1,161 of its 4,718 strains in any of 397 runs, so the
full-scale run is a query change and is the largest single change available to the figure.
And the graph-prior probe decides whether nine attention heads are enforcing a premise
nobody has checked. Both land before any arm ladder is worth launching;
\cref{sec:campaign} sizes what follows.

\notesec{Figure 6} The betaxanthin case is the strongest thing in this document and it does
not depend on beating anyone on accuracy. The argument is coverage: yeast-GEM reaches
$19\,\%$ of the screen and misses $73\,\%$ of its high producers, and CGT finds high
producers in the part a flux method cannot represent at $4.4\times$ over base rate.

The last panel is where the figure either becomes an argument or stays a list, and getting
its \emph{structure} right is the open work. The sequence that carries it: setup diagram;
betaxanthin performance against its $0.914$ ceiling; the coverage gap; the capability gap;
the Flux Cone Learning comparison, framed as a tie on the genes both can score; and then
the structural panel. That last one is the constraint-based layer above: the amino-acid
coupling is real in the data and unused by a plain auxiliary head, and a stoichiometric
constraint over yeast-GEM is what would close it. Shown that way the amino-acid material is
the \emph{motivation} for the mechanism rather than a failed joint-training result, which
is the only framing in which a negative belongs in a figure.

Keep the option text in the draw.io boxes. Several of those panels name alternatives that
have not been run, and a composition that still reads when one of them does not arrive is
worth more than a tidier one that does not.
